\documentclass[12pt, a4paper]{article}

% --- PACOTES PRINCIPAIS ---
\usepackage[utf8]{inputenc}
\usepackage[portuguese]{babel} % Ajustado para português
\usepackage[T1]{fontenc}
\usepackage{amsmath, amsfonts, amssymb}
\usepackage{geometry}
\usepackage{booktabs}
\usepackage{array}
\usepackage{graphicx}
\usepackage{caption}
\usepackage{float}
\usepackage{siunitx}
\usepackage{hyperref}

% --- CONFIGURAÇÃO DAS MARGENS (ABNT NBR 14724) ---
\geometry{top=3cm, bottom=2cm, left=3cm, right=2cm}

% --- CONFIGURAÇÕES DO SIUNITX (ISO 80000-1) ---
\sisetup{
  output-decimal-marker = {,},
  separate-uncertainty = true,
  per-mode = symbol
}

\linespread{1.5}

\begin{document}

% --- CABEÇALHO E CAPA ---
\begin{center}
    \textbf{\Large UNIVERSIDADE DE BRASÍLIA -- UnB}\\
    \textbf{\large Faculdade do Gama -- FGA}\\
    \vspace{0.5cm}
    \textbf{\large Disciplina: Física 1 Experimental}\\
    \vspace{1.5cm}
    {\Large \textbf{RELATÓRIO DE EXPERIMENTO}}\\
    {\large \textbf{Experimento II - Movimento Retilíneo Uniforme (MRU)}}\\
    \vspace{1.5cm}
\end{center}

\begin{flushright}
    \textbf{Integrantes do Grupo:}\\
    Emanuel Matsumori Quintiliano -- Matrícula: 261040090\\
    Felipe Wermelinger da Costa Leite -- Matrícula: 261014993\\
    Rodrigo Silveira Pereira Costa -- Matrícula: 251027826\\
    Yago dos Santos Rezende -- Matrícula: 241012427\\
    \vspace{0.5cm}
    \textbf{Professor:} Prof. Dr. Wytler C. Santos\\
    \textbf{Data do Experimento:} 27 de setembro de 2026
\end{flushright}

\vfill
\begin{center}
    Brasília -- DF \\
    2026
\end{center}

\newpage

% --- OBJETIVOS ---
\section{Objetivos}
% TODO: Colegas, preencham/melhorem os objetivos baseando-se no roteiro.
Este experimento visa descrever a cinemática de um movimento retilíneo e uniforme através do acompanhamento da posição e da velocidade com a evolução do tempo. Pretende-se introduzir e estudar a propagação de erros e a aplicação de métodos gráficos para a análise cinemática, utilizando um trilho de ar com atrito desprezável.

% --- INTRODUÇÃO TEÓRICA ---
\section{Introdução Teórica}
% TODO: Colegas, expliquem a teoria do MRU (Movimento Retilíneo Uniforme) aqui. 
A equação horária que descreve o movimento retilíneo uniforme (MRU) é dada por:
\begin{equation}
    x(t) = x_0 + v \cdot t
\end{equation}
Onde $x(t)$ representa a posição no instante $t$, $x_0$ a posição inicial e $v$ a velocidade constante do móvel. Esta secção deve também explicar como a regressão linear permite extrair os parâmetros físicos ($v$ e $x_0$) a partir dos coeficientes angular e linear do gráfico.

% --- PARTE EXPERIMENTAL E DISCUSSÃO ---
\section{Parte Experimental e Resultados}

\subsection{Materiais Utilizados}
O equipamento base utilizado consistiu num kit de trilho de ar. Os principais componentes foram:
\begin{itemize}
    \item 01 trilho de 120 cm;
    \item 01 cronómetro digital multifunções com fonte DC 12 V;
    \item 02 sensores fotoelétricos com suporte fixador (S1 e S2);
    \item 01 eletroíman com bornes e haste;
    \item 01 chave liga-desliga;
    % TODO: Colegas, adicionem mais materiais relevantes se acharem necessário.
\end{itemize}

\subsection{Procedimentos}
% TODO: Colegas, descrevam resumidamente os passos da medição (ajuste dos sensores, as massas penduradas, as repetições das contagens no cronómetro, etc).
As medições consistiram em posicionar o primeiro sensor ($S_1$) na posição $x_0 = 0\text{ cm}$ e variar a distância do segundo sensor ($S_2$) em intervalos de $\SI{10}{\cm}$ até atingir a distância de $\SI{50}{\cm}$. O cronómetro foi accionado pela passagem do carrinho pelos sensores e, para cada deslocamento, foram efetuadas 5 medições de tempo.

\subsection{Dados Experimentais}
A Tabela~\ref{tab:dados} apresenta os valores dos intervalos de tempo medidos, o tempo médio e os respetivos erros calculados para as distâncias pré-definidas.

\begin{table}[!htbp]
\centering
\caption{Valores experimentais para o Movimento Retilíneo Uniforme.}
\label{tab:dados}
\vspace{0.2cm}
\resizebox{\textwidth}{!}{%
\begin{tabular}{lccccc}
\toprule
\textbf{Parâmetro} & $\Delta x_1 = \SI{10}{\cm}$ & $\Delta x_2 = \SI{20}{\cm}$ & $\Delta x_3 = \SI{30}{\cm}$ & $\Delta x_4 = \SI{40}{\cm}$ & $\Delta x_5 = \SI{50}{\cm}$ \\
\midrule
$\Delta t_1$ (s) & 0,237 & 0,438 & 0,690 & 0,882 & 1,255 \\
$\Delta t_2$ (s) & 0,233 & 0,470 & 0,686 & 0,913 & 1,251 \\
$\Delta t_3$ (s) & 0,236 & 0,464 & 0,694 & 0,900 & 1,248 \\
$\Delta t_4$ (s) & 0,267 & 0,465 & 0,686 & 0,899 & 1,238 \\
$\Delta t_5$ (s) & 0,265 & 0,468 & 0,702 & 0,919 & 1,236 \\
\midrule
\textbf{Tempo médio (s)} & 0,248 & 0,461 & 0,691 & 0,902 & 1,246 \\
\textbf{Erro aleatório (s)} & 0,007 & 0,005 & 0,003 & 0,000 & 0,003 \\
\textbf{Erro absoluto (s)} & 0,008 & 0,006 & 0,004 & 0,007 & 0,004 \\
\bottomrule
\end{tabular}%
}
\end{table}

\subsection{Análise Gráfica e Regressão Linear}
Para extrair a velocidade constante do movimento, procedeu-se ao ajuste linear dos dados de posição em função do tempo médio. A equação teórica da reta apresenta a forma $y = A \cdot x + B$, em que $y$ corresponde ao deslocamento ($\Delta x$) e $x$ ao tempo médio ($\Delta t$).

A análise computacional obteve os seguintes coeficientes para a regressão linear:
\begin{itemize}
    \item \textbf{Coeficiente Angular (Declive):} $A = 40,62 \pm 2,37 \text{ cm/s}$
    \item \textbf{Coeficiente Linear (Ordenada na origem):} $B = 1,18 \pm 1,87 \text{ cm}$
\end{itemize}

Comparando a equação da reta ($y = Ax + B$) com a equação horária da cinemática ($x = v \cdot t + x_0$), conclui-se que o coeficiente angular $A$ representa a velocidade média do móvel ($v \approx 40,62 \text{ cm/s}$) e o coeficiente linear $B$ representa a posição inicial ajustada ($x_0 \approx 1,18 \text{ cm}$). O valor de $B$ está muito próximo de zero, o que corrobora o facto de o referencial ter sido ajustado para partir de $x_0 = 0$.

% TODO: Adicionar o gráfico gerado computacionalmente (substituir o nome do ficheiro)
\begin{figure}[H]
    \centering
    %\includegraphics[width=0.8\textwidth]{nome_da_imagem_do_grafico.png}
    \caption{Gráfico da regressão linear relacionando o deslocamento e o tempo.}
    \label{fig:grafico}
\end{figure}

% --- CONCLUSÃO ---
\section{Conclusão}
% TODO: Colegas, redijam a conclusão final aqui. 
% Discutam se a equação horária obtida condiz com as expectativas e se a aproximação de que v = A e x_0 = B se sustenta.

% --- REFERÊNCIAS ---
\section*{Referências}
\begin{enumerate}
    \item SANTOS, Wytler C. \textbf{Roteiro para o Experimento II - Movimento Retilíneo Uniforme (MRU)}. Faculdade UnB Gama.
\end{enumerate}

\end{document}
